\section{Conclusion}
\label{sec:conclusion}

We have applied the NestSection algorithm to all 49 bicameral state legislatures,
producing the first comprehensive empirical study of nested bicameral redistricting
at national scale.

The main findings are:
\begin{enumerate}
  \item \textbf{Compatibility}: 41 of 49 bicameral states have
    $\gcd(H,S) > 1$ and are compatible with NestSection. The remaining
    8 are incompatible due to coprime chamber sizes; changing one
    chamber's size (a constitutional amendment in most cases) would achieve
    compatibility.
  \item \textbf{Compactness penalty}: The nesting constraint costs $+2.5\%$
    in mean edge-cut (house: $+2.1\%$, senate: $+3.4\%$) and
    $-0.008$/$-0.028$ PP points for house/senate respectively.
    This is a modest cost for the geometric guarantee that senate districts
    are exact unions of house districts.
  \item \textbf{Partisan neutrality}: Nesting changes partisan outcomes by
    at most $0.3$ seats on average. The nesting constraint is not a partisan
    device; its effect on partisan outcomes is negligible.
  \item \textbf{Population balance}: All 41 compatible states achieve
    $\leq 0.5\%$ maximum deviation in both chambers simultaneously.
\end{enumerate}

These findings position NestSection as a practical tool for state redistricting
commissions in the 41 compatible states. The algorithm is implemented in the
\texttt{redist} Rust binary via \texttt{redist state --chamber nest} and
runs in under 10 minutes per state at block-group resolution.

\paragraph{Constitutional implications.}
In states with constitutional or statutory nesting requirements, NestSection
provides a computationally efficient mechanism for satisfying those requirements
while optimising compactness. Our results suggest that the legal requirement
to nest senate districts within house districts does not materially compromise
either compactness or partisan fairness---the two primary substantive concerns
in state legislative redistricting litigation.

\paragraph{Future work.}
The next step is to apply NestSection to the 8 incompatible states using
ad-hoc nesting procedures and to compare outcomes. A complementary analysis
would examine multi-level nesting: states where county legislative districts
must nest within state legislative districts, which nest within congressional
districts. This three-level nesting problem requires a multi-spine variant of
NestSection that is the subject of ongoing research.
